\chapter{Laços}

\section{\texttt{for}}

Itera sobre uma sequência. A variável de iteração é visível dentro do bloco e
\textit{persiste} após o laço com o último valor assumido:

\begin{lstlisting}[caption={for sobre range}]
for i in 1..4 {
  display i    // 1, 2, 3  (4 não incluído)
}
display i    // 3 — variável persiste após o laço
\end{lstlisting}

\subsection*{Range exclusivo e inclusivo}

\hayashi{} segue a mesma convenção de Rust:

\begin{lstlisting}
for i in 1..3  { }    // 1, 2      (exclusivo: não inclui o fim)
for i in 1..=3 { }    // 1, 2, 3   (inclusivo: inclui o fim)
\end{lstlisting}

Ranges também são expressões que avaliam para listas — ver Seção~\ref{sec:ranges}.

\subsection*{Iterando sobre listas}

\begin{lstlisting}
let frutas = ["maçã", "banana", "uva"]

for fruta in frutas {
  display fruta
}
\end{lstlisting}

\subsection*{Descartando a variável de iteração}

Quando o índice não importa:

\begin{lstlisting}
for _ in 1..3 {
  display "repetição"
}
\end{lstlisting}

\section{\texttt{while}}

Executa enquanto a condição for verdadeira:

\begin{lstlisting}[caption={while básico}]
let n = 1

while n <= 5 {
  display n
  n += 1
}
\end{lstlisting}

\begin{aviso}
  \hayashi{} não detecta laços infinitos. Um \lstinline|while true { }| sem
  \lstinline{break} executará indefinidamente. Use \texttt{Ctrl-C} para interromper.
\end{aviso}

\section{\texttt{break} e \texttt{continue}}

\lstinline{break} interrompe o laço imediatamente. \lstinline{continue} avança para
a próxima iteração:

\begin{lstlisting}[caption={break e continue}]
for i in 1..10 {
  if i == 3 { continue }    // pula o 3
  if i == 7 { break }       // para no 7
  display i
}
// exibe: 1, 2, 4, 5, 6
\end{lstlisting}

\section{Laços com DataFrames}

\lstinline{for} pode iterar sobre listas produzidas por funções de dados:

\begin{lstlisting}[caption={Iterando sobre grupos}]
load "vendas.csv" as df

for uf in ["RS", "SP", "RJ"] {
  let sub = filter(df, estado == uf)
  display f"{uf}: {count(sub)} obs"
}
\end{lstlisting}
